Goodman and Kruskal's $\gamma$ \citeyearpar{goodman-kruskal-1954} measures association between two ordinal variables. It compares the numbers of concordant pairs of observations, $C$, and discordant pairs, $D$. A concordant pair has the same ordering on both variables; a discordant pair has opposite orderings. The coefficient is

\begin{equation}
\gamma^{GK} = \frac{C-D}{C+D}.
\end{equation}

Partial $\gamma$ \citep{davis-1967} measures association within strata defined by the conditioning variables. For example, testing DIF between $Y_i$ and $X_j$ requires conditioning on the total score $S$. We therefore stratify the data by the observed values of $S$ and count concordant and discordant pairs within each stratum. Let $n$ denote the number of strata, indexed by $l=1,\dots,n$ to keep the stratum index distinct from the item indices $i,j$. Pooling the counts across strata gives the partial gamma coefficient, which can also be expressed as a weighted sum of the stratum-specific gamma coefficients:
\begin{align}
    \hat{\gamma}_{Y_i\bot X_j\mid S} &= \sum_{l=1}^n w_l\gamma^{GK}_l = \sum_{l=1}^n\frac{(C_l+D_l)}{\sum_{l=1}^n(C_l+D_l)}\frac{C_l-D_l}{C_l+D_l} \\
    &= \frac{\sum_{l=1}^n(C_l-D_l)}{\sum_{l=1}^n(C_l+D_l)} = \frac{\sum_{l=1}^nC_l-\sum_{l=1}^nD_l}{\sum_{l=1}^nC_l+\sum_{l=1}^nD_l}.
\end{align}
Each stratum's weight is its number of comparable pairs (concordant plus discordant) divided by the total number of comparable pairs across all strata. Strata with more comparable pairs therefore contribute more to the pooled coefficient \citep{agresti-1984}. If a stratum $l$ has no comparable pairs, $C_l+D_l=0$, its stratum-specific coefficient $\gamma^{GK}_l$ is undefined and the weighted-sum expression contains the indeterminate term $0\cdot(0/0)$. We adopt the convention that such a stratum contributes nothing to either sum, i.e., it is simply omitted from the stratification: since $C_l=D_l=0$ leaves both $\sum_l C_l$ and $\sum_l D_l$ unchanged, this is exactly what the pooled-count form on the right already does, and it matches how Algorithm~\ref{alg:part_gam} in Appendix~\ref{sec:AppAlgorithms} is computed, which skips any stratum whose contingency table has fewer than two rows or columns.

Both coefficients compare the difference between concordant and discordant pair counts with their sum. When defined, they lie in $[-1,1]$: values closer to $1$ indicate stronger positive association, and values closer to $-1$ indicate stronger negative association \citep{agresti-1984}. If there are no comparable pairs, the denominator is zero and the coefficient is undefined.
